% ==================================================================================
% Combined Baseline Architectures — Condensed for NeurIPS Manuscript
% Include via: % ==================================================================================
% Combined Baseline Architectures — Condensed for NeurIPS Manuscript
% Include via: % ==================================================================================
% Combined Baseline Architectures — Condensed for NeurIPS Manuscript
% Include via: % ==================================================================================
% Combined Baseline Architectures — Condensed for NeurIPS Manuscript
% Include via: \input{combined_baselines.tex}
% Assumes parent document defines: \mR, \op, \mat, and loads amsmath, bm, algorithm, algpseudocode
% ==================================================================================

\subsection{Baseline Architectures}
\label{sec:baselines}

All four baselines share a common infrastructure: a FasterRCNN-ResNet50 backbone for monocular 3D detection, an LKS (Last Known State) memory buffer for non-differentiable visual persistence, a GlobalStructuralEncoder that converts oriented bounding box (OBB) corners into per-object 3D structural tokens, an InterObjectTransformer for intra-frame spatial reasoning with continuous pairwise 3D positional encodings, and a three-head relationship predictor (attention, spatial, contacting). They differ in which optional temporal and spatial modules are activated, forming a controlled ablation ladder. Table~\ref{tab:baseline_modules} summarizes the module activation pattern.

\begin{table}[h]
\centering
\caption{Module activation across baselines. \checkmark\,= active, --\,= absent.}
\label{tab:baseline_modules}
\small
\begin{tabular}{lcccc}
\toprule
\textbf{Module} & \textbf{W-STTran} & \textbf{W-STTran++} & \textbf{W-DsgDetr} & \textbf{W-DsgDetr++} \\
\midrule
LKS Buffer + GlobalStructuralEncoder  & \checkmark & \checkmark & \checkmark & \checkmark \\
ObjectSpatialEncoder                   & --         & \checkmark & --         & \checkmark \\
ObjectMotionEncoder + MotionFusion     & --         & \checkmark & --         & \checkmark \\
CameraTemporalEncoder                  & --         & --         & --         & \checkmark \\
TemporalObjectEncoder                  & --         & --         & \checkmark & \checkmark \\
TemporalEdgeAttention                  & --         & \checkmark & \checkmark & \checkmark \\
\bottomrule
\end{tabular}
\end{table}

% ----------------------------------------------------------------------------------
\paragraph{W-STTran.}
\label{par:wsttran}
W-STTran is the minimal baseline, adapting STTran~to the WorldSGG task. Object features extracted by the FasterRCNN backbone are held in the LKS buffer, which returns the nearest visible frame's feature $\mat{m}_n^{(t)}$ for occluded objects and a staleness signal $\Delta_n^{(t)}$. The GlobalStructuralEncoder converts each object's 8-corner OBB $\mat{C}_n \in \mR^{8\times3}$ into a centered, translation-invariant structural token $\mat{g}_n$ via a shared MLP. The LKS Tokenizer fuses geometry, buffered features, and log-staleness through a linear projection with LayerNorm: $\mat{x}_n^{(t)} = \op{FusionProj}([\mat{g}_n^{(t)} \;\|\; \mat{m}_n^{(t)} \;\|\; \log(\Delta_n^{(t)}+1)])$. The InterObjectTransformer then performs global spatial reasoning using additive pairwise positional encodings derived from 3D inter-object distances and volume ratios. A Node Predictor classifies objects, and the Relationship Predictor forms edge tokens from subject/object features, union RoI features, and CLIP semantic embeddings, refines them via self-attention, and outputs predictions through three independent MLP heads. W-STTran uses \emph{no} camera, motion, or temporal modules, isolating the contribution of passive 3D memory alone.

% ----------------------------------------------------------------------------------
\paragraph{W-STTran++.}
\label{par:wsttranpp}
W-STTran++ augments the minimal baseline with camera-relative and motion-aware representations, but still omits explicit temporal object tracking. The ObjectSpatialEncoder computes per-object camera-relative features from the extrinsic matrix $\mat{T}=[\mat{R}\,|\,\boldsymbol{\tau}]$: log-distance, view alignment $\alpha_n = \hat{\mat{d}}_n \cdot (-\mat{R}_{:,2})$, and azimuth components, projected via MLP to $\mat{c}_n \in \mR^{d_{\text{camera}}}$. The ObjectMotionEncoder captures 3D dynamics by computing world-frame velocity $\mat{v}_n^{(t)}$ and acceleration $\mat{a}_n^{(t)}$ from OBB center finite differences, alongside camera-relative velocity $\mat{R}^{\!\top}\mat{v}_n$, yielding an 11-dimensional feature projected to $\mat{\mu}_n^{(t)}$. The Tokenizer concatenates geometry, buffer, spatial, and staleness features; the resulting token is then fused with motion via $\hat{\mat{x}}_n^{(t)} = \op{LayerNorm}(\op{Linear}([\mat{x}_n^{(t)} \;\|\; \mat{\mu}_n^{(t)}]))$. After InterObjectTransformer reasoning, TemporalEdgeAttention performs cross-frame self-attention over each unique object pair's relationship tokens with learned temporal positional embeddings, enabling temporal consistency without explicit per-object tracking.

% ----------------------------------------------------------------------------------
\paragraph{W-DsgDetr.}
\label{par:wdsgdetr}
W-DsgDetr adapts DSG-DETR~to WorldSGG by introducing explicit per-object temporal tracking. It shares the LKS buffer, GlobalStructuralEncoder, and Tokenizer with W-STTran (without camera or motion features). The key addition is the TemporalObjectEncoder: for each object $n$, it collects the $T$-length sequence of tokenized representations and processes them with a Transformer encoder using sinusoidal temporal positional encodings: $\tilde{\mat{X}}_n = \op{TransformerEncoder}(\{\mat{x}_n^{(t)} + \mat{PE}_{\text{sin}}(t)\}_{t=1}^T)$. This allows the model to interpolate and contextualize object features across time, addressing occlusion gaps that the zero-order LKS buffer cannot. The temporally-refined tokens then undergo InterObjectTransformer spatial reasoning, followed by node prediction, relationship formation with intra-frame self-attention, and TemporalEdgeAttention for cross-frame edge consistency. This design isolates the contribution of explicit temporal tracking relative to W-STTran.

% ----------------------------------------------------------------------------------
\paragraph{W-DsgDetr++.}
\label{par:wdsgdetrpp}
W-DsgDetr++ activates every available module, representing the most feature-rich baseline before WorldWise. Beyond W-DsgDetr's TemporalObjectEncoder, it adds the ObjectSpatialEncoder, ObjectMotionEncoder, and a CameraTemporalEncoder that models ego-motion dynamics. The CameraTemporalEncoder computes relative camera poses between consecutive frames ($\mat{R}_{\text{rel}}^{(t)} = \mat{R}_t\mat{R}_{t-1}^{\!\top}$, $\boldsymbol{\tau}_{\text{rel}}^{(t)} = \boldsymbol{\tau}_t - \mat{R}_{\text{rel}}^{(t)}\boldsymbol{\tau}_{t-1}$) and encodes the sequence via self-attention to produce per-frame ego-motion tokens $\mat{e}_t$. The MotionFusion stage integrates object motion and ego-motion: $\hat{\mat{x}}_n^{(t)} = \op{LayerNorm}(\op{Linear}([\mat{x}_n^{(t)} \;\|\; \mat{\mu}_n^{(t)} \;\|\; \mat{e}_t]))$. The full pipeline then proceeds through TemporalObjectEncoder, InterObjectTransformer, and TemporalEdgeAttention. This configuration tests the upper bound of performance achievable within the passive-memory (LKS) paradigm.

% ----------------------------------------------------------------------------------
\paragraph{Training Objective.}
\label{par:baseline_loss}
All four baselines share a common loss function that partitions human-object pairs into \emph{visible} (both entities observed) and \emph{unseen} (at least one occluded) buckets. Visible pairs use standard cross-entropy (attention) and binary cross-entropy (spatial, contacting) with clean ground truth. Unseen pairs employ VLM pseudo-labels with label smoothing and a weighting factor $\lambda_{\text{vlm}}$:
\begin{equation}
    \mathcal{L} = \sum_{r \in \{\text{att, spa, con}\}} \bigl(\mathcal{L}_{r}^{\text{vis}} + \mathcal{L}_{r}^{\text{unseen}}\bigr) + \mathcal{L}_{\text{node}}.
    \label{eq:baseline_loss}
\end{equation}

% Assumes parent document defines: \mR, \op, \mat, and loads amsmath, bm, algorithm, algpseudocode
% ==================================================================================

\subsection{Baseline Architectures}
\label{sec:baselines}

All four baselines share a common infrastructure: a FasterRCNN-ResNet50 backbone for monocular 3D detection, an LKS (Last Known State) memory buffer for non-differentiable visual persistence, a GlobalStructuralEncoder that converts oriented bounding box (OBB) corners into per-object 3D structural tokens, an InterObjectTransformer for intra-frame spatial reasoning with continuous pairwise 3D positional encodings, and a three-head relationship predictor (attention, spatial, contacting). They differ in which optional temporal and spatial modules are activated, forming a controlled ablation ladder. Table~\ref{tab:baseline_modules} summarizes the module activation pattern.

\begin{table}[h]
\centering
\caption{Module activation across baselines. \checkmark\,= active, --\,= absent.}
\label{tab:baseline_modules}
\small
\begin{tabular}{lcccc}
\toprule
\textbf{Module} & \textbf{W-STTran} & \textbf{W-STTran++} & \textbf{W-DsgDetr} & \textbf{W-DsgDetr++} \\
\midrule
LKS Buffer + GlobalStructuralEncoder  & \checkmark & \checkmark & \checkmark & \checkmark \\
ObjectSpatialEncoder                   & --         & \checkmark & --         & \checkmark \\
ObjectMotionEncoder + MotionFusion     & --         & \checkmark & --         & \checkmark \\
CameraTemporalEncoder                  & --         & --         & --         & \checkmark \\
TemporalObjectEncoder                  & --         & --         & \checkmark & \checkmark \\
TemporalEdgeAttention                  & --         & \checkmark & \checkmark & \checkmark \\
\bottomrule
\end{tabular}
\end{table}

% ----------------------------------------------------------------------------------
\paragraph{W-STTran.}
\label{par:wsttran}
W-STTran is the minimal baseline, adapting STTran~to the WorldSGG task. Object features extracted by the FasterRCNN backbone are held in the LKS buffer, which returns the nearest visible frame's feature $\mat{m}_n^{(t)}$ for occluded objects and a staleness signal $\Delta_n^{(t)}$. The GlobalStructuralEncoder converts each object's 8-corner OBB $\mat{C}_n \in \mR^{8\times3}$ into a centered, translation-invariant structural token $\mat{g}_n$ via a shared MLP. The LKS Tokenizer fuses geometry, buffered features, and log-staleness through a linear projection with LayerNorm: $\mat{x}_n^{(t)} = \op{FusionProj}([\mat{g}_n^{(t)} \;\|\; \mat{m}_n^{(t)} \;\|\; \log(\Delta_n^{(t)}+1)])$. The InterObjectTransformer then performs global spatial reasoning using additive pairwise positional encodings derived from 3D inter-object distances and volume ratios. A Node Predictor classifies objects, and the Relationship Predictor forms edge tokens from subject/object features, union RoI features, and CLIP semantic embeddings, refines them via self-attention, and outputs predictions through three independent MLP heads. W-STTran uses \emph{no} camera, motion, or temporal modules, isolating the contribution of passive 3D memory alone.

% ----------------------------------------------------------------------------------
\paragraph{W-STTran++.}
\label{par:wsttranpp}
W-STTran++ augments the minimal baseline with camera-relative and motion-aware representations, but still omits explicit temporal object tracking. The ObjectSpatialEncoder computes per-object camera-relative features from the extrinsic matrix $\mat{T}=[\mat{R}\,|\,\boldsymbol{\tau}]$: log-distance, view alignment $\alpha_n = \hat{\mat{d}}_n \cdot (-\mat{R}_{:,2})$, and azimuth components, projected via MLP to $\mat{c}_n \in \mR^{d_{\text{camera}}}$. The ObjectMotionEncoder captures 3D dynamics by computing world-frame velocity $\mat{v}_n^{(t)}$ and acceleration $\mat{a}_n^{(t)}$ from OBB center finite differences, alongside camera-relative velocity $\mat{R}^{\!\top}\mat{v}_n$, yielding an 11-dimensional feature projected to $\mat{\mu}_n^{(t)}$. The Tokenizer concatenates geometry, buffer, spatial, and staleness features; the resulting token is then fused with motion via $\hat{\mat{x}}_n^{(t)} = \op{LayerNorm}(\op{Linear}([\mat{x}_n^{(t)} \;\|\; \mat{\mu}_n^{(t)}]))$. After InterObjectTransformer reasoning, TemporalEdgeAttention performs cross-frame self-attention over each unique object pair's relationship tokens with learned temporal positional embeddings, enabling temporal consistency without explicit per-object tracking.

% ----------------------------------------------------------------------------------
\paragraph{W-DsgDetr.}
\label{par:wdsgdetr}
W-DsgDetr adapts DSG-DETR~to WorldSGG by introducing explicit per-object temporal tracking. It shares the LKS buffer, GlobalStructuralEncoder, and Tokenizer with W-STTran (without camera or motion features). The key addition is the TemporalObjectEncoder: for each object $n$, it collects the $T$-length sequence of tokenized representations and processes them with a Transformer encoder using sinusoidal temporal positional encodings: $\tilde{\mat{X}}_n = \op{TransformerEncoder}(\{\mat{x}_n^{(t)} + \mat{PE}_{\text{sin}}(t)\}_{t=1}^T)$. This allows the model to interpolate and contextualize object features across time, addressing occlusion gaps that the zero-order LKS buffer cannot. The temporally-refined tokens then undergo InterObjectTransformer spatial reasoning, followed by node prediction, relationship formation with intra-frame self-attention, and TemporalEdgeAttention for cross-frame edge consistency. This design isolates the contribution of explicit temporal tracking relative to W-STTran.

% ----------------------------------------------------------------------------------
\paragraph{W-DsgDetr++.}
\label{par:wdsgdetrpp}
W-DsgDetr++ activates every available module, representing the most feature-rich baseline before WorldWise. Beyond W-DsgDetr's TemporalObjectEncoder, it adds the ObjectSpatialEncoder, ObjectMotionEncoder, and a CameraTemporalEncoder that models ego-motion dynamics. The CameraTemporalEncoder computes relative camera poses between consecutive frames ($\mat{R}_{\text{rel}}^{(t)} = \mat{R}_t\mat{R}_{t-1}^{\!\top}$, $\boldsymbol{\tau}_{\text{rel}}^{(t)} = \boldsymbol{\tau}_t - \mat{R}_{\text{rel}}^{(t)}\boldsymbol{\tau}_{t-1}$) and encodes the sequence via self-attention to produce per-frame ego-motion tokens $\mat{e}_t$. The MotionFusion stage integrates object motion and ego-motion: $\hat{\mat{x}}_n^{(t)} = \op{LayerNorm}(\op{Linear}([\mat{x}_n^{(t)} \;\|\; \mat{\mu}_n^{(t)} \;\|\; \mat{e}_t]))$. The full pipeline then proceeds through TemporalObjectEncoder, InterObjectTransformer, and TemporalEdgeAttention. This configuration tests the upper bound of performance achievable within the passive-memory (LKS) paradigm.

% ----------------------------------------------------------------------------------
\paragraph{Training Objective.}
\label{par:baseline_loss}
All four baselines share a common loss function that partitions human-object pairs into \emph{visible} (both entities observed) and \emph{unseen} (at least one occluded) buckets. Visible pairs use standard cross-entropy (attention) and binary cross-entropy (spatial, contacting) with clean ground truth. Unseen pairs employ VLM pseudo-labels with label smoothing and a weighting factor $\lambda_{\text{vlm}}$:
\begin{equation}
    \mathcal{L} = \sum_{r \in \{\text{att, spa, con}\}} \bigl(\mathcal{L}_{r}^{\text{vis}} + \mathcal{L}_{r}^{\text{unseen}}\bigr) + \mathcal{L}_{\text{node}}.
    \label{eq:baseline_loss}
\end{equation}

% Assumes parent document defines: \mR, \op, \mat, and loads amsmath, bm, algorithm, algpseudocode
% ==================================================================================

\subsection{Baseline Architectures}
\label{sec:baselines}

All four baselines share a common infrastructure: a FasterRCNN-ResNet50 backbone for monocular 3D detection, an LKS (Last Known State) memory buffer for non-differentiable visual persistence, a GlobalStructuralEncoder that converts oriented bounding box (OBB) corners into per-object 3D structural tokens, an InterObjectTransformer for intra-frame spatial reasoning with continuous pairwise 3D positional encodings, and a three-head relationship predictor (attention, spatial, contacting). They differ in which optional temporal and spatial modules are activated, forming a controlled ablation ladder. Table~\ref{tab:baseline_modules} summarizes the module activation pattern.

\begin{table}[h]
\centering
\caption{Module activation across baselines. \checkmark\,= active, --\,= absent.}
\label{tab:baseline_modules}
\small
\begin{tabular}{lcccc}
\toprule
\textbf{Module} & \textbf{W-STTran} & \textbf{W-STTran++} & \textbf{W-DsgDetr} & \textbf{W-DsgDetr++} \\
\midrule
LKS Buffer + GlobalStructuralEncoder  & \checkmark & \checkmark & \checkmark & \checkmark \\
ObjectSpatialEncoder                   & --         & \checkmark & --         & \checkmark \\
ObjectMotionEncoder + MotionFusion     & --         & \checkmark & --         & \checkmark \\
CameraTemporalEncoder                  & --         & --         & --         & \checkmark \\
TemporalObjectEncoder                  & --         & --         & \checkmark & \checkmark \\
TemporalEdgeAttention                  & --         & \checkmark & \checkmark & \checkmark \\
\bottomrule
\end{tabular}
\end{table}

% ----------------------------------------------------------------------------------
\paragraph{W-STTran.}
\label{par:wsttran}
W-STTran is the minimal baseline, adapting STTran~to the WorldSGG task. Object features extracted by the FasterRCNN backbone are held in the LKS buffer, which returns the nearest visible frame's feature $\mat{m}_n^{(t)}$ for occluded objects and a staleness signal $\Delta_n^{(t)}$. The GlobalStructuralEncoder converts each object's 8-corner OBB $\mat{C}_n \in \mR^{8\times3}$ into a centered, translation-invariant structural token $\mat{g}_n$ via a shared MLP. The LKS Tokenizer fuses geometry, buffered features, and log-staleness through a linear projection with LayerNorm: $\mat{x}_n^{(t)} = \op{FusionProj}([\mat{g}_n^{(t)} \;\|\; \mat{m}_n^{(t)} \;\|\; \log(\Delta_n^{(t)}+1)])$. The InterObjectTransformer then performs global spatial reasoning using additive pairwise positional encodings derived from 3D inter-object distances and volume ratios. A Node Predictor classifies objects, and the Relationship Predictor forms edge tokens from subject/object features, union RoI features, and CLIP semantic embeddings, refines them via self-attention, and outputs predictions through three independent MLP heads. W-STTran uses \emph{no} camera, motion, or temporal modules, isolating the contribution of passive 3D memory alone.

% ----------------------------------------------------------------------------------
\paragraph{W-STTran++.}
\label{par:wsttranpp}
W-STTran++ augments the minimal baseline with camera-relative and motion-aware representations, but still omits explicit temporal object tracking. The ObjectSpatialEncoder computes per-object camera-relative features from the extrinsic matrix $\mat{T}=[\mat{R}\,|\,\boldsymbol{\tau}]$: log-distance, view alignment $\alpha_n = \hat{\mat{d}}_n \cdot (-\mat{R}_{:,2})$, and azimuth components, projected via MLP to $\mat{c}_n \in \mR^{d_{\text{camera}}}$. The ObjectMotionEncoder captures 3D dynamics by computing world-frame velocity $\mat{v}_n^{(t)}$ and acceleration $\mat{a}_n^{(t)}$ from OBB center finite differences, alongside camera-relative velocity $\mat{R}^{\!\top}\mat{v}_n$, yielding an 11-dimensional feature projected to $\mat{\mu}_n^{(t)}$. The Tokenizer concatenates geometry, buffer, spatial, and staleness features; the resulting token is then fused with motion via $\hat{\mat{x}}_n^{(t)} = \op{LayerNorm}(\op{Linear}([\mat{x}_n^{(t)} \;\|\; \mat{\mu}_n^{(t)}]))$. After InterObjectTransformer reasoning, TemporalEdgeAttention performs cross-frame self-attention over each unique object pair's relationship tokens with learned temporal positional embeddings, enabling temporal consistency without explicit per-object tracking.

% ----------------------------------------------------------------------------------
\paragraph{W-DsgDetr.}
\label{par:wdsgdetr}
W-DsgDetr adapts DSG-DETR~to WorldSGG by introducing explicit per-object temporal tracking. It shares the LKS buffer, GlobalStructuralEncoder, and Tokenizer with W-STTran (without camera or motion features). The key addition is the TemporalObjectEncoder: for each object $n$, it collects the $T$-length sequence of tokenized representations and processes them with a Transformer encoder using sinusoidal temporal positional encodings: $\tilde{\mat{X}}_n = \op{TransformerEncoder}(\{\mat{x}_n^{(t)} + \mat{PE}_{\text{sin}}(t)\}_{t=1}^T)$. This allows the model to interpolate and contextualize object features across time, addressing occlusion gaps that the zero-order LKS buffer cannot. The temporally-refined tokens then undergo InterObjectTransformer spatial reasoning, followed by node prediction, relationship formation with intra-frame self-attention, and TemporalEdgeAttention for cross-frame edge consistency. This design isolates the contribution of explicit temporal tracking relative to W-STTran.

% ----------------------------------------------------------------------------------
\paragraph{W-DsgDetr++.}
\label{par:wdsgdetrpp}
W-DsgDetr++ activates every available module, representing the most feature-rich baseline before WorldWise. Beyond W-DsgDetr's TemporalObjectEncoder, it adds the ObjectSpatialEncoder, ObjectMotionEncoder, and a CameraTemporalEncoder that models ego-motion dynamics. The CameraTemporalEncoder computes relative camera poses between consecutive frames ($\mat{R}_{\text{rel}}^{(t)} = \mat{R}_t\mat{R}_{t-1}^{\!\top}$, $\boldsymbol{\tau}_{\text{rel}}^{(t)} = \boldsymbol{\tau}_t - \mat{R}_{\text{rel}}^{(t)}\boldsymbol{\tau}_{t-1}$) and encodes the sequence via self-attention to produce per-frame ego-motion tokens $\mat{e}_t$. The MotionFusion stage integrates object motion and ego-motion: $\hat{\mat{x}}_n^{(t)} = \op{LayerNorm}(\op{Linear}([\mat{x}_n^{(t)} \;\|\; \mat{\mu}_n^{(t)} \;\|\; \mat{e}_t]))$. The full pipeline then proceeds through TemporalObjectEncoder, InterObjectTransformer, and TemporalEdgeAttention. This configuration tests the upper bound of performance achievable within the passive-memory (LKS) paradigm.

% ----------------------------------------------------------------------------------
\paragraph{Training Objective.}
\label{par:baseline_loss}
All four baselines share a common loss function that partitions human-object pairs into \emph{visible} (both entities observed) and \emph{unseen} (at least one occluded) buckets. Visible pairs use standard cross-entropy (attention) and binary cross-entropy (spatial, contacting) with clean ground truth. Unseen pairs employ VLM pseudo-labels with label smoothing and a weighting factor $\lambda_{\text{vlm}}$:
\begin{equation}
    \mathcal{L} = \sum_{r \in \{\text{att, spa, con}\}} \bigl(\mathcal{L}_{r}^{\text{vis}} + \mathcal{L}_{r}^{\text{unseen}}\bigr) + \mathcal{L}_{\text{node}}.
    \label{eq:baseline_loss}
\end{equation}

% Assumes parent document defines: \mR, \op, \mat, and loads amsmath, bm, algorithm, algpseudocode
% ==================================================================================

\subsection{Baseline Architectures}
\label{sec:baselines}

All four baselines share a common infrastructure: a FasterRCNN-ResNet50 backbone for monocular 3D detection, an LKS (Last Known State) memory buffer for non-differentiable visual persistence, a GlobalStructuralEncoder that converts oriented bounding box (OBB) corners into per-object 3D structural tokens, an InterObjectTransformer for intra-frame spatial reasoning with continuous pairwise 3D positional encodings, and a three-head relationship predictor (attention, spatial, contacting). They differ in which optional temporal and spatial modules are activated, forming a controlled ablation ladder. Table~\ref{tab:baseline_modules} summarizes the module activation pattern.

\begin{table}[h]
\centering
\caption{Module activation across baselines. \checkmark\,= active, --\,= absent.}
\label{tab:baseline_modules}
\small
\begin{tabular}{lcccc}
\toprule
\textbf{Module} & \textbf{W-STTran} & \textbf{W-STTran++} & \textbf{W-DsgDetr} & \textbf{W-DsgDetr++} \\
\midrule
LKS Buffer + GlobalStructuralEncoder  & \checkmark & \checkmark & \checkmark & \checkmark \\
ObjectSpatialEncoder                   & --         & \checkmark & --         & \checkmark \\
ObjectMotionEncoder + MotionFusion     & --         & \checkmark & --         & \checkmark \\
CameraTemporalEncoder                  & --         & --         & --         & \checkmark \\
TemporalObjectEncoder                  & --         & --         & \checkmark & \checkmark \\
TemporalEdgeAttention                  & --         & \checkmark & \checkmark & \checkmark \\
\bottomrule
\end{tabular}
\end{table}

% ----------------------------------------------------------------------------------
\paragraph{W-STTran.}
\label{par:wsttran}
W-STTran is the minimal baseline, adapting STTran~to the WorldSGG task. Object features extracted by the FasterRCNN backbone are held in the LKS buffer, which returns the nearest visible frame's feature $\mat{m}_n^{(t)}$ for occluded objects and a staleness signal $\Delta_n^{(t)}$. The GlobalStructuralEncoder converts each object's 8-corner OBB $\mat{C}_n \in \mR^{8\times3}$ into a centered, translation-invariant structural token $\mat{g}_n$ via a shared MLP. The LKS Tokenizer fuses geometry, buffered features, and log-staleness through a linear projection with LayerNorm: $\mat{x}_n^{(t)} = \op{FusionProj}([\mat{g}_n^{(t)} \;\|\; \mat{m}_n^{(t)} \;\|\; \log(\Delta_n^{(t)}+1)])$. The InterObjectTransformer then performs global spatial reasoning using additive pairwise positional encodings derived from 3D inter-object distances and volume ratios. A Node Predictor classifies objects, and the Relationship Predictor forms edge tokens from subject/object features, union RoI features, and CLIP semantic embeddings, refines them via self-attention, and outputs predictions through three independent MLP heads. W-STTran uses \emph{no} camera, motion, or temporal modules, isolating the contribution of passive 3D memory alone.

% ----------------------------------------------------------------------------------
\paragraph{W-STTran++.}
\label{par:wsttranpp}
W-STTran++ augments the minimal baseline with camera-relative and motion-aware representations, but still omits explicit temporal object tracking. The ObjectSpatialEncoder computes per-object camera-relative features from the extrinsic matrix $\mat{T}=[\mat{R}\,|\,\boldsymbol{\tau}]$: log-distance, view alignment $\alpha_n = \hat{\mat{d}}_n \cdot (-\mat{R}_{:,2})$, and azimuth components, projected via MLP to $\mat{c}_n \in \mR^{d_{\text{camera}}}$. The ObjectMotionEncoder captures 3D dynamics by computing world-frame velocity $\mat{v}_n^{(t)}$ and acceleration $\mat{a}_n^{(t)}$ from OBB center finite differences, alongside camera-relative velocity $\mat{R}^{\!\top}\mat{v}_n$, yielding an 11-dimensional feature projected to $\mat{\mu}_n^{(t)}$. The Tokenizer concatenates geometry, buffer, spatial, and staleness features; the resulting token is then fused with motion via $\hat{\mat{x}}_n^{(t)} = \op{LayerNorm}(\op{Linear}([\mat{x}_n^{(t)} \;\|\; \mat{\mu}_n^{(t)}]))$. After InterObjectTransformer reasoning, TemporalEdgeAttention performs cross-frame self-attention over each unique object pair's relationship tokens with learned temporal positional embeddings, enabling temporal consistency without explicit per-object tracking.

% ----------------------------------------------------------------------------------
\paragraph{W-DsgDetr.}
\label{par:wdsgdetr}
W-DsgDetr adapts DSG-DETR~to WorldSGG by introducing explicit per-object temporal tracking. It shares the LKS buffer, GlobalStructuralEncoder, and Tokenizer with W-STTran (without camera or motion features). The key addition is the TemporalObjectEncoder: for each object $n$, it collects the $T$-length sequence of tokenized representations and processes them with a Transformer encoder using sinusoidal temporal positional encodings: $\tilde{\mat{X}}_n = \op{TransformerEncoder}(\{\mat{x}_n^{(t)} + \mat{PE}_{\text{sin}}(t)\}_{t=1}^T)$. This allows the model to interpolate and contextualize object features across time, addressing occlusion gaps that the zero-order LKS buffer cannot. The temporally-refined tokens then undergo InterObjectTransformer spatial reasoning, followed by node prediction, relationship formation with intra-frame self-attention, and TemporalEdgeAttention for cross-frame edge consistency. This design isolates the contribution of explicit temporal tracking relative to W-STTran.

% ----------------------------------------------------------------------------------
\paragraph{W-DsgDetr++.}
\label{par:wdsgdetrpp}
W-DsgDetr++ activates every available module, representing the most feature-rich baseline before WorldWise. Beyond W-DsgDetr's TemporalObjectEncoder, it adds the ObjectSpatialEncoder, ObjectMotionEncoder, and a CameraTemporalEncoder that models ego-motion dynamics. The CameraTemporalEncoder computes relative camera poses between consecutive frames ($\mat{R}_{\text{rel}}^{(t)} = \mat{R}_t\mat{R}_{t-1}^{\!\top}$, $\boldsymbol{\tau}_{\text{rel}}^{(t)} = \boldsymbol{\tau}_t - \mat{R}_{\text{rel}}^{(t)}\boldsymbol{\tau}_{t-1}$) and encodes the sequence via self-attention to produce per-frame ego-motion tokens $\mat{e}_t$. The MotionFusion stage integrates object motion and ego-motion: $\hat{\mat{x}}_n^{(t)} = \op{LayerNorm}(\op{Linear}([\mat{x}_n^{(t)} \;\|\; \mat{\mu}_n^{(t)} \;\|\; \mat{e}_t]))$. The full pipeline then proceeds through TemporalObjectEncoder, InterObjectTransformer, and TemporalEdgeAttention. This configuration tests the upper bound of performance achievable within the passive-memory (LKS) paradigm.

% ----------------------------------------------------------------------------------
\paragraph{Training Objective.}
\label{par:baseline_loss}
All four baselines share a common loss function that partitions human-object pairs into \emph{visible} (both entities observed) and \emph{unseen} (at least one occluded) buckets. Visible pairs use standard cross-entropy (attention) and binary cross-entropy (spatial, contacting) with clean ground truth. Unseen pairs employ VLM pseudo-labels with label smoothing and a weighting factor $\lambda_{\text{vlm}}$:
\begin{equation}
    \mathcal{L} = \sum_{r \in \{\text{att, spa, con}\}} \bigl(\mathcal{L}_{r}^{\text{vis}} + \mathcal{L}_{r}^{\text{unseen}}\bigr) + \mathcal{L}_{\text{node}}.
    \label{eq:baseline_loss}
\end{equation}
